% Generated by `python -m experiments.experiment_1.appendix_captions` -- do not edit by hand.
% Methods/appendix text for the main-text figure fig:exp1-main (main_figure.tex).
% Place it under a heading carrying the label below, e.g.
% \subsection{Design-space figure: details}\label{app:exp1-main-details}
Each random design in the top panel of Figure~\ref{fig:exp1-main} has two regressors, each continuous or binary (one contrast, so $u_k(:c_1)$ denotes a slope or a contrast), and their interaction, with each parameter entering each coefficient at random; the set of 12 designs was redrawn until together they cover every coefficient of the grid. Each design contributes 200 simulated datasets of 500 trials, with 200 posterior draws per dataset. Coefficients are shown on the unconstrained (link) scale on which the network is trained; $\pi$ (middle panel) is shown on its natural scale, with prior $\mathrm{Beta}(1, 50)$. For the design comparison (bottom panel), both DDM predictors act on drift rate and boundary separation (11 estimated coefficients, including $\pi$). Expected information gain is the Gaussian log-volume ratio $\tfrac{1}{2}\log\det\Sigma_{\text{prior}} - \tfrac{1}{2}\,\mathbb{E}_x\!\left[\log\det\Sigma_{\text{post}}(x)\right]$, in nats, with $\Sigma_{\text{prior}}$ estimated from 1500 prior draws; NRMSE is the root mean squared error of the posterior draws against ground truth, divided by the standard deviation of the ground-truth values and averaged over the coefficients. Both are computed on a common estimand, each slope scaled by the realised standard deviation of its predictor, so that continuous and binary predictors are comparable. Information gain is relative to each design's own prior, which is wider for binary contrasts (a contrast of two cell means has twice the prior variance of a continuous slope). Each point uses 200 datasets with 200 posterior draws; standard errors are over datasets (bootstrapped for NRMSE). Binary predictors are randomly unbalanced dichotomisations (realised balance 0.38--0.45), not median splits.
